\chapter[Introdução]{Introdução}

A competição Micromouse, anunciada pela revista IEEE Spectrum na década de 1970 \cite{wikipedia_micromouse}, desafia robôs autônomos de pequeno porte a solucionarem labirintos desconhecidos no menor tempo possível. Para realizar essa 
tarefa de forma independente, esses robôs utilizam princípios fundamentais de navegação, como mapeamento, planejamento de trajetória e localização frequentemente apoiados em métodos de busca como o algoritmo de Flood Fill \cite{wikipedia_micromouse}.
A busca por maior eficiência na exploração de labirintos tem sido foco de diversas pesquisas científicas. Estudos recentes propõem otimizações matemáticas na navegação, como a modificação do algoritmo Flood Fill para considerar o
custo temporal das rotações do robô em seu próprio eixo \cite{cai2012}. Outras abordagens incluem o algoritmo ET-Floodfill, que introduz uma taxa de custa esperada, chamada expected toll, ao lidar com paredes desconhecidas, tratando matematicamente
a incerteza do ambiente para tornar o planejamento de rotas mais inteligente \cite{cai2012}.

No cenário mundial, a robótica móvel desponta como uma ferramenta com enorme potencial educacional para cativar estudantes nas áreas de STEM (Ciência, Tecnologia, Engenharia e Matemática), utilizando competições como o Micromouse para estimular a
aplicação prática de conhecimentos multidisciplinares \cite{zawadniak2020}. Enquanto competidores de altíssimo desempenho desenvolvem robôs complexos, equipados até mesmo com ventoinhas geradores de vácuo para garantir maior aderência e aceleração na pista
\cite{wikipedia_micromouse}, existe paralelalmente uma forte necessidade de sistemas acessíveis e didáticos para as universidades. Em resposta ao alto custo das tecnologias de ponta, organizações internacionais têm promovido o desenvolvimento de plataformas
robóticas de entrada construídas com materiais e componentes baratos, visando viabilizar a introdução de novos estudantes no meio competitivo \cite{ukmars}. Dessa forma, constata-se que o mercado-alvo e os maiores consumidores finais para esses sistemas englobam
a comunidade acadêmica, professores, pesquisadores e entusiastas que buscam plataformas tangíveis e acessíveis para o ensino integrado de engenharias.

Sob esse contexto, o desenvolvimento de um projeto voltado à competição Micromouse configura-se como um valioso campo de testes práticos, exigindo dos participantes não apenas domínio teórico, mas habilidades concretas de execução técnica. Este projeto justifica-se
pela necessidade inerente de preparar profissionais capacitados para lidarem com sistemas interligados complexos, unindo simultaneamente hardware, software e estruturas mecânicas em um único ecossistema aplicado ao mundo real. A elaboração de um robô inteligente, 
capaz de operar ativamente em conjunto com um sistema web para exibição de telemetria em tempo real, atende à alta demanda das universidades por plataformas educacionais robustas. O produto final melhora a abordagem de ensino tradicional ao superar as limitações
do aprendizado estritamente teórico, entregando uma solução tecnológica completa, autônoma e multidisciplinar.